\documentclass[10pt,letterpaper]{article}
\usepackage{cornell-notes}

\title{Annex A.04.04: 04-04-structuring-formalisms-and-structural-categories}
\author{}
\date{}
\setDocTitle{Annex A.04.04: 04-04-structuring-formalisms-and-structural-categories}
\setDocSubtitle{ISO/IEC/IEEE 42010:2022}
\setDocOwner{ISO 42010 Cornell Notes}
\setDocDate{\today}
\setCornellCollection{ISO/IEC/IEEE 42010:2022 Cornell Notes}
\setCornellUnitType{Annex}
\setCornellUnitNumber{A.04.04}
\setCornellUnitTitle{04-04-structuring-formalisms-and-structural-categories}
\hypersetup{pdftitle={Annex A.04.04: 04-04-structuring-formalisms-and-structural-categories},pdfsubject={Cornell notes},pdfauthor={}}

\begin{document}
\maketitle
\begin{CornellOverviewBox}{Overview}
{\Large\bfseries\textcolor{CNavy}{Annex A.4.4: Structuring formalisms and structural categories}}\\[3pt]
{\small\textbf{Classification:} Informative annex}\\
{\small\textbf{Learning objective:} Use the informative guidance on structuring formalisms and structural categories to reinforce the normative and conceptual material without treating the guidance itself as mandatory.}
\end{CornellOverviewBox}
\vspace{3pt}

\begin{CornellNotesTable}
\CornellNoteRow{Purpose / key concept}{An ADF can provide a structuring formalism, i.e. a set of rules for using architecture considerations and correspondences between them, to organize the architecture viewpoints used to generate associated views, e.g. a grid framework formalism.}
\CornellNoteRow{Core details}{\begin{itemize}[leftmargin=*,nosep,topsep=1pt]
\item The purpose of the structuring formalism is to provide ways of representing relationships among various elements of the architecture and enhancing opportunities for analysis of interactions among those elements.
\item Here, this concept of framework dimensions is called “structural categories” since the way these categories are used in various ADFs is not always aligned with the concept of “dimensions.” So, the structural categories term is used herein since it is a broader and more inclusive concept.
\item An ADF can define structural categories used in the structuring formalism.
\item Sometimes these categories are represented by “dimensions” in a graphic portrayal, such as the rows and columns used in several frameworks.
\end{itemize}}
\CornellNoteRow{Example / application}{Well known structuring formalisms include: GERAM cube in ISO 15704, Reference Architectural Model Industrie 4.0 (RAMI 4.0), TOGAF stack (business, information systems, technology), NAF grid, UAF grid, and Zachman Framework matrix.}
\CornellNoteRow{Related sections}{4.0}
\CornellNoteRow{Active recall}{\begin{itemize}[leftmargin=*,nosep,topsep=1pt]
\item What is the central role of Structuring formalisms and structural categories?
\item How does Structuring formalisms and structural categories connect to adjacent architecture-description concepts?
\end{itemize}}
\end{CornellNotesTable}


\begin{CornellSummaryBox}{Summary}
An ADF can provide a structuring formalism, i.e. a set of rules for using architecture considerations and correspondences between them, to organize the architecture viewpoints used to generate associated views, e.g. a grid framework formalism. This material is informative guidance and does not itself create a conformance requirement.
\end{CornellSummaryBox}

\end{document}
